\documentclass[10pt,a4paper]{amsart}
\usepackage[margin=19mm]{geometry}
\usepackage{amsmath,amssymb,mathtools}
\usepackage[scr=rsfs,cal=euler]{mathalfa}
\usepackage{xfrac}
\usepackage[unicode,hidelinks,bookmarks=false]{hyperref}
\newcommand{\ie}{i.e.\ }
\newcommand{\cf}{cf.\ }
\newcommand{\resp}{respectively\ }
\newcommand{\Subsection}{\subsection}
\providecommand{\nobreakdash}{\nobreak\hbox{-}}
\setlength{\emergencystretch}{2em}
\multlinegap0pt
\usepackage{bm}
\usepackage{bbm}
\usepackage{dsfont}
\usepackage{xspace}
\usepackage{cancel}
\usepackage{mathtools}
\usepackage{amsthm}
\makeatletter
\@ifundefined{lemma}{\newtheorem{lemma}{Lemma}}{}
\@ifundefined{theorem}{\newtheorem{theorem}[lemma]{Theorem}}{}
\makeatother
\newcommand{\Deg}{\operatorname{Deg}}
\newcommand{\R}{\mathbb{R}}
\newcommand{\dom}{\textnormal{dom}} % domain
\title[On the accretivity and m-accretivity of Laplacians and porou]{On the accretivity and m-accretivity of Laplacians and porous medium-type operators on graphs}
\author{Davide Bianchi, Matthias Keller, Alberto G. Setti, Rados{\l}aw K. Wojciechowski}
\date{Source: arXiv:2607.09625v2 (CC BY 4.0)}
\begin{document}
\maketitle

\noindent\emph{Source and licence.} On the accretivity and m-accretivity of Laplacians and porous medium-type operators on graphs. Version arXiv:2607.09625v2 (CC BY 4.0), distributed under the Creative Commons Attribution 4.0 International licence, as recorded in the arXiv licence metadata for this exact version and shown on the abstract page https://arxiv.org/abs/2607.09625v2. Full rights notice: licenses/arxiv-2607.09625-CC-BY-4.0.txt.\par\smallskip

\noindent\textit{Editorial scope.} Counted unit: the Theorem whose source label is \texttt{thm:min} in the article (source0.tex lines 738--774, source label \texttt{thm:min}), delivered together with the article's own complete original proof of it (source0.tex lines 775--985, one \texttt{proof} environment of 211 lines). No mathematics is written, repaired or re-derived: statement and proof are the article's own text line for line. Required context retained uncounted: thm:min\_principle (lines 621--647); eq:ineq\_1 (lines 667--670). Every definition, display and label of those lines is retained unchanged. Omitted from this package and not redistributed: 1--620, 648--666, 671--737, 986--4419. Nothing inside a retained excerpt is omitted or altered: every retained body is delivered exactly as the article's lines are written, so no omission receipt of any kind accompanies this package. Citations: the retained excerpts carry no citation command at all, so no bibliography entry of the article is transcribed and no reference is redistributed. Reference adaptation: none was needed and none was made; every reference inside the delivered unit resolves inside it, so no unresolved reference or citation is produced. Printed slips kept literal: no printed word, symbol, hypothesis or number of the counted statement or of its proof was altered. Layout and class: the article's own class is \texttt{article} with packages including arxiv, babel, csquotes, booktabs, graphicx, sidecap, hyperref, amssymb, amsmath, amsthm, mathrsfs, mathtools, accents, ulem; the wrapper is the lane's pinned \texttt{amsart} editorial wrapper at 10pt on A4, which restates the theorem-like environments the retained text needs and adds the article's own macro definitions that the retained text needs (\texttt{\textbackslash Deg}, \texttt{\textbackslash R}, \texttt{\textbackslash dom}), with extra wrapper packages (bm, bbm, dsfont, xspace, cancel, mathtools). The delivered numbering is the unit's own. Assets: no retained span contains a figure environment, a graphics-inclusion command, a PDF-page inclusion, a table or a TikZ picture; no image file is copied into this package, the package declares no asset dependency and no image file is redistributed with it. Licence: version arXiv:2607.09625v2 of this article is distributed under the Creative Commons Attribution 4.0 International licence, as recorded in the arXiv licence metadata for this exact version and shown on the abstract page https://arxiv.org/abs/2607.09625v2. A full rights notice is delivered at licenses/arxiv-2607.09625-CC-BY-4.0.txt. Characters: every retained excerpt is pure ASCII, so the delivered engine, XeLaTeX with its default Latin Modern text font, reports no missing character.
\begin{theorem}[Sign preservation]\label{thm:min_principle}
Let $\lambda>0$, $\sigma \colon X \to (0, \infty)$ and $\phi, \,\psi \colon \R \to \R$ be monotone increasing, with
$\phi(0)=0$ and $\psi(t)<0$ for $t<0$.
Let $u \in \dom(\Delta \Phi)$ satisfy 
$$\left( \sigma \Psi + \lambda\Delta \Phi\right)u \geq 0.$$ 
We consider four cases:

\begin{itemize}
	\item[\emph{\textbf{Case 1:}}]
	There exists $x_0 \in X$ such that $u$ attains a minimum at $x_0$, i.e.,
	$$u(x_{0})=\inf_{x\in X} u(x) > -\infty.$$
	
	\item[\emph{\textbf{Case 2:}}] $G$ does not contain any infinite paths. 	

    \item[\emph{\textbf{Case 3:}}] For every infinite path $(x_{n})$, we have $\sum_{n} \mu(x_{n})=\infty$ and there exists  a monotonically increasing  function $\gamma \colon [0,\infty) \to [0,\infty)$ with  $\gamma(t)>0$ for $t>0$ such  that 
    $$
    \sum_{n}\gamma \bigl( |u(x_{n})| \bigr) \mu(x_{n})<\infty .
    $$
  \item[\emph{\textbf{Case 4:}}] The function $u$ is bounded below 
  and, for every infinite path $(x_{n})$, we have 
    $$\sum_{n} \frac{\sigma(x_n)}{\Deg(x_{n})}=\infty.$$
\end{itemize}    
	 In all four cases, $u \geq 0$. 
     
     Moreover, in all cases,
     { assuming additionally that $0=\phi(0)<\phi(t)$ for all $t>0$ and $\psi(0)\leq 0$}, if $u(x)=0$ for some $x\in X$, then  $u= 0$. 
\end{theorem}

\begin{equation}\label{eq:ineq_1}
\sum_{y \in X} w(x_0,y) \Phi u(y) \leq { \Deg(x_0)\mu(x_0)
\left[\frac{\sigma(x_0)}{\lambda \Deg(x_0)} \Psi u(x_0) +  \Phi u(x_0) \right].}
\end{equation}

\begin{theorem}[Comparison principle]\label{thm:min}
Let $\lambda>0$, $\sigma \colon X \to (0, \infty)$ and { $\phi\colon\R\to\R$ be continuous and monotone increasing with $\phi(0)=0$}. 
Let $u_k \in \dom(\Delta \Phi)$ and  $g_k \in C(X)$ satisfy
	$$
	\left(\sigma + \lambda\Delta \Phi \right) u_k = g_k \qquad \mbox{ for } k=1,2.
	$$
{ Set
\[
 u\coloneqq u_1-u_2
 \qquad \text{and} \qquad
 v\coloneqq \Phi u_1-\Phi u_2.
\]
If $g_1\geq g_2$ and $G$ and $v$ satisfy the assumptions of one of the Cases~1--3 of
Theorem~\ref{thm:min_principle}, then
$u_1\geq u_2$. The same conclusion holds in Case~4 provided that, for 
every $c>0$, there exists $\varepsilon_c>0$ such that
\begin{equation}\label{extra_assumption_case4}
 v(x)\leq-c
 \quad\Longrightarrow\quad
 u(x)\leq-\varepsilon_c \qquad \text{for all } x \in X.
\end{equation}} 

In particular, if  $u_1$ and  $u_2$ are in $\ell^p(X,\mu)$ for some $p \in [1,\infty)$, then { $v$ satisfies the summability condition in
Case~3}. 
If $u_1$ and $u_2$ are in $\ell^\infty(X)$, { then $v$ is bounded and satisfies
the additional condition \eqref{extra_assumption_case4}}.
It follows that 
\begin{itemize}
\item[\textup{(i)}] If $\sum_n \mu(x_n)=\infty$ for every infinite path $(x_n)$, then
$\sigma + \lambda\Delta \Phi$ is injective on $\ell^p(X,\mu) \cap \dom(\Delta \Phi)$
for all $p\in[1,\infty)$ { and its inverse is order preserving on its range}. 
\item[\textup{(ii)}] If $\sum_{n} \sigma(x_n)/\Deg(x_{n})=\infty$ for every infinite path $(x_n)$,
then $\sigma + \lambda\Delta \Phi$ is injective on
{ $\ell^\infty(X)\cap\dom(\Delta\Phi)$ and its inverse is order preserving on its range}.
\end{itemize}

\end{theorem}

\begin{proof}
{
Notice that   $v\in\dom(\Delta)$, and subtracting the  equations satisfied by $u_1$ and $u_2$ gives
\begin{equation}
\label{eq:comparison-differences}
 \sigma u+\lambda\Delta v=g_1-g_2\geq0. 
\end{equation}
Moreover, monotonicity of $\phi$ implies
\begin{equation}
\label{eq:weak-sign-compatibility}
 \{v<0\}\subseteq\{u<0\}\subseteq\{v\leq0\}. 
\end{equation}

We now adapt the descent step in the proof of Theorem~\ref{thm:min_principle}
to show
that if $x_0\in X$ satisfies $u(x_0)<0$, then there exists
$x_1\sim x_0$ such that $v(x_1)<v(x_0) \leq 0$.
Indeed, let $x_0\in X$ satisfy
$u(x_0)<0$, so that $v(x_0)\leq0$ by \eqref{eq:weak-sign-compatibility}. If $x_0$ has no neighbors, then
\[
 \Delta v(x_0)=\frac{\kappa(x_0)}{\mu(x_0)}v(x_0)\leq0
\]
which contradicts \eqref{eq:comparison-differences} at $x_0$ since $u(x_0)<0$.
Thus, $x_0$ has at least one neighbor and $\Deg(x_0)>0$. Rearranging
\eqref{eq:comparison-differences}, we obtain
\begin{equation}
\label{eq:ineq_1bis}
 \sum_{y\in X}w(x_0,y)v(y)
 \leq
 \Deg(x_0) \mu(x_0)\left[ \frac{\sigma(x_0)}{\lambda\Deg(x_0)}u(x_0)
 +v(x_0)\right]. 
\end{equation}
The same weighted-average argument used after \eqref{eq:ineq_1} in the
proof of Theorem~\ref{thm:min_principle}, with $\Phi u$ and $\Psi u$
there replaced by $v$ and $u$, applied to \eqref{eq:ineq_1bis} shows that
there exists $x_1\sim x_0$ such that
\begin{equation}
\label{eq:v-ineq}
 v(x_1)
 \leq v(x_0)+\frac{\sigma(x_0)}{\lambda\Deg(x_0)}u(x_0)=:A_{x_0}
 <v(x_0)\leq0. 
\end{equation}
In particular, $v(x_1)<0$ and therefore $u(x_1)<0$ by
\eqref{eq:weak-sign-compatibility}, so the descent step can be
iterated.

We now show that $u \geq 0$. Suppose, by contradiction, that $u(z)<0$ at some vertex $z$. By the preceding
descent step, there exists $x_0\sim z$ such that
\[
v(x_0)<v(z)\leq0.
\]
Thus, $v(x_0)<0$ and $u(x_0)<0$ by
\eqref{eq:weak-sign-compatibility}. In
Case~1, applying~\eqref{eq:v-ineq} at a point where
$v$ attains its negative minimum gives a still smaller value. In
Case~2, iteration gives an infinite path of descent, since $v$ strictly decreases
and hence all its vertices are distinct. In Case~3, the resulting
infinite path $(x_n)$ satisfies $v(x_{n+1})<v(x_n)<0$, and therefore
\[
 \sum_n\gamma\bigl(|v(x_n)|\bigr)\mu(x_n)
 \geq
 \gamma\bigl(|v(x_0)|\bigr)\sum_n\mu(x_n)=\infty
\]
which gives a contradiction.

In Case~4, let $c=-v(x_0)>0$. Along the descending path we have
$v(x_n)\leq-c$, so~\eqref{extra_assumption_case4} gives
$u(x_n)\leq-\varepsilon_c$. Iterating
\eqref{eq:v-ineq} yields
\[
 v(x_{n+1})
 \leq v(x_0)-\frac{\varepsilon_c}{\lambda}
 \sum_{k=0}^n\frac{\sigma(x_k)}{\Deg(x_k)}
 \to -\infty
\]
as $n \to \infty$, contradicting the fact that $v$ is bounded below.
Thus, the assumption $u(z)<0$ leads to a contradiction in Cases~1--3
and also in Case~4 under \eqref{extra_assumption_case4}. Hence, $u\geq0$.
}
{ Thus, we have shown $u_1\geq u_2$ in Cases~1--3 and in
Case~4 under condition~\eqref{extra_assumption_case4}.

We now assume that $u_1,u_2 \in \ell^p(X,\mu)$ for some $p \in [1,\infty)$
and show that $v$ satisfies the summability condition in Case 3 and that (i) holds.}
{
Define $\alpha\colon[0,\infty)\to[0,\infty)$ by
\[
 \alpha(t):=\phi(t)-\phi(-t).
\]
The assumptions on $\phi$ imply directly that $\alpha$ is continuous,
positive and monotone increasing, with $\alpha(0)=0$.
We define now the generalized  inverse by
\[
 \beta(s):=\inf\{t\geq0 \mid \alpha(t)\geq s\}
 \qquad \text{for } s\geq0,
\]
where $\inf\emptyset:=\infty$, and let
\[
 \gamma(s):=\min\{1,\beta(s)^p\}.
\]
Here and below, $\min\{1,\infty\}:=1$. Both $\beta$ and $\gamma$ are
monotone increasing and $\beta(0)=\gamma(0)=0$.

For $s>0$, continuity of $\alpha$ at zero gives the existence of $\delta_s>0$ such that
$\alpha(t)<s$ for $0\leq t<\delta_s$. Hence, $\beta(s)\geq\delta_s$ if
the defining set for $\beta(s)$ is nonempty, whereas $\beta(s)=\infty$
otherwise. Thus, $\gamma(s)>0$ for every $s>0$. Moreover, $t$ is
admissible in the definition of $\beta(\alpha(t))$, so
$\beta(\alpha(t))\leq t$ and, hence,
\begin{equation}\label{eq:gamma-alpha-estimate}
 \gamma(\alpha(t))
 =
 \min\{1,\beta(\alpha(t))^p\}
 \leq\min\{1,t^p\}
 \leq t^p.
\end{equation}

Now, for every $x\in X$, let
\[
 M(x):=\max\{|u_1(x)|,|u_2(x)|\}.
\]
Both $u_1(x)$ and $u_2(x)$ belong to $[-M(x),M(x)]$. Since $\phi$ is
monotone increasing, both $\phi(u_1(x))$ and $\phi(u_2(x))$ belong to
\[
 [\phi(-M(x)),\phi(M(x))].
\]
It follows that
\begin{align*}
 |v(x)|
 &=
 |\phi(u_1(x))-\phi(u_2(x))|\leq
 \phi(M(x))-\phi(-M(x))
 =
 \alpha(M(x)).
\end{align*}
Since $\gamma$ is monotone increasing, we obtain from
\eqref{eq:gamma-alpha-estimate} that
\[
 \gamma(|v(x)|)
 \leq
 \gamma(\alpha(M(x)))
 \leq
 M(x)^p.
\]
Finally,
\[
 M(x)^p
 =
 \max\{|u_1(x)|^p,|u_2(x)|^p\}
 \leq
 |u_1(x)|^p+|u_2(x)|^p.
\]
Therefore, for every infinite path $(x_n)$,
\begin{align*}
 \sum_n\gamma(|v(x_n)|)\mu(x_n)
 \leq
 \|u_1\|_p^p+\|u_2\|_p^p
 <\infty.
\end{align*}
Thus, $v$ satisfies the summability condition in Case~3 of
Theorem~\ref{thm:min_principle}. If every infinite path has infinite
measure, then all of the hypotheses of Case~3 are satisfied.
{ The statements of (i) now follow from what we have proved above.

We now assume $u_1,u_2 \in \ell^\infty(X)$ and show $v \in \ell^\infty(X)$, 
that $u$ and $v$ satisfy \eqref{extra_assumption_case4} and (ii) holds.}
Let
\[
 M:=\max\{\|u_1\|_\infty,\|u_2\|_\infty\}.
\]
Then, $u_1(x),u_2(x)\in[-M,M]$ for every $x\in X$ and $v$ is
bounded since $\phi$ is continuous on $[-M,M]$.

Fix $c>0$. By the uniform continuity of $\phi$ on $[-M,M]$, there exists
$\varepsilon_c>0$ such that
\[
 |r-s|<\varepsilon_c
 \quad\Longrightarrow\quad
 |\phi(r)-\phi(s)|<c
 \qquad\text{for all }r,s\in[-M,M].
\]
Equivalently, by contraposition,
\[
 |\phi(r)-\phi(s)|\geq c
 \quad\Longrightarrow\quad
 |r-s|\geq\varepsilon_c.
\]
If $v(x)\leq-c$, then
\[
 \phi(u_2(x))-\phi(u_1(x))=-v(x)\geq c>0.
\]
Since $\phi$ is monotone increasing, this also implies $u_2(x)>u_1(x)$.
Consequently,
\[
 u_2(x)-u_1(x)
 =
 |u_2(x)-u_1(x)|
 \geq\varepsilon_c
\]
and, hence,
\[
 u(x)=u_1(x)-u_2(x)\leq-\varepsilon_c.
\]
Thus, the additional condition~\eqref{extra_assumption_case4} is
satisfied.

The preceding comparison gives order preservation on the range.
Applying it in both directions when $g_1=g_2$ gives injectivity.}
{ This proves the statements in (ii) and completes the proof.}

\end{proof}

\end{document}
